\chapter{Introduction}
\label{ch:introduction}

\section{Background and Motivation}
\label{sec:intro-background}

Motion in modern sport is fast and highly variable, which makes consistent
manual measurement difficult. As recording has spread from broadcast crews to
ordinary phones and consumer cameras, computer vision (CV) has become a
practical tool for sports analytics, applied to problems from player and ball
tracking through to tactical and performance analysis
\cite{thomas2017cvsports}. Tennis is a representative and demanding case:
matches are decided by the trajectory of a small, rapidly moving ball over a
precisely measured court, so an accurate spatial and temporal understanding of
both the ball and the playing surface is fundamental to almost any downstream
analysis.

The commercial prominence of vision-based officiating illustrates both the
value and the difficulty of this problem. Systems such as Hawk-Eye are now
embedded in professional tennis for automated line-calling, yet scholarly
analysis has cautioned that their confident, deterministic visualisations can
obscure the measurement uncertainty inherent in any reconstruction of ball
position \cite{collins2008hawkeye}. Such systems also depend on calibrated
multi-camera installations that are unavailable outside elite venues. This
motivates a complementary line of work: vision-only analysis of ordinary
monocular video that is reproducible, inexpensive, and accessible for coaching,
amateur play, and research.

Automated tennis analysis from a single video stream rests on three perception
sub-problems. The first is ball tracking: localising the ball in every
frame. This is a canonical small-object problem, because the ball occupies only
a handful of pixels, moves quickly enough to blur or leave afterimages, and is
frequently occluded by players, the net, or court markings, so that
single-frame appearance is often insufficient. Heatmap-based temporal models
such as TrackNet were introduced precisely to exploit motion cues across
consecutive frames in this setting \cite{huang2019tracknet}, in contrast to
single-frame object detectors of the YOLO family that localise objects through
bounding-box regression \cite{redmon2016yolo}. The second sub-problem is
court keypoint detection and registration: recovering the geometric
mapping between image pixels and real-world court coordinates, typically by
detecting a fixed set of court landmarks and fitting a homography with a robust
estimator such as random sample consensus (RANSAC) \cite{fischler1981ransac}.
Keypoint localisation has been tackled both with high-capacity backbones such
as deep residual networks \cite{he2016resnet} and high-resolution pose networks
\cite{sun2019hrnet}, and with efficient mobile architectures designed for
constrained hardware \cite{howard2019mobilenetv3}. The third sub-problem is
bounce-event detection: identifying the discrete frames at which the ball
contacts the court, recovered from a noisy and intermittently missing
trajectory. This is a temporal pattern-recognition task that can be addressed
with hand-crafted kinematic heuristics, gradient-boosted decision trees
\cite{chen2016xgboost}, or temporal convolutional networks (TCNs)
\cite{bai2018tcn}.

A wide range of models has been proposed for each of these sub-problems
individually. However, they are commonly reported on heterogeneous datasets and
private train/test splits, and most studies evaluate a single proposed model
rather than placing competing designs on a common footing. As a result, it is
difficult to draw like-for-like conclusions about the accuracy, speed, and
robustness trade-offs between architectures within a sub-task, and, to
the best of the author's knowledge, comparatively few works compose the
resulting components into a single, end-to-end pipeline whose practical
deployability can be assessed. The present thesis is motivated by both gaps: the scarcity of controlled
within-task comparison, and the still rarer attempt to assemble the resulting
components into one reproducible pipeline.

\section{Problem Statement}
\label{sec:intro-problem}

This thesis studies the three perception sub-tasks of monocular tennis video
analysis (ball tracking, court keypoint detection with homography, and
bounce-event detection), together with the engineering problem of integrating
them into a single pipeline. Each sub-task is framed as a controlled model
comparison conducted under a shared, deterministic evaluation protocol.

The sub-tasks are formulated as follows. Ball tracking is the per-frame
estimation of the ball's image position, with an explicit ``not visible'' state
for frames in which the ball is absent or undetectable. Court detection is the
per-frame localisation of a fixed set of fourteen predefined court keypoints,
followed by estimation of a homography that maps the image onto a canonical
court template expressed in metric units. Bounce detection is the
classification of which frames correspond to a ball bounce, distinguished both
from racket-hit frames and from the much larger set of ordinary in-flight
frames.

A central methodological point is that this work draws on two distinct
datasets. Ball tracking and bounce detection share the publicly available
TrackNet tennis dataset, in which each frame is annotated with ball visibility,
image coordinates, and a status label. Court keypoint detection uses a separate
court-keypoint dataset with its own annotation schema and image resolution.
Because the two datasets differ in content, resolution, and split policy,
accuracy figures are comparable only within a sub-task and never across
sub-tasks; this constraint is stated explicitly and respected throughout the
thesis. The integration sub-problem is distinct from the three perception
problems: composing the strongest model from each sub-task into one system that
ingests raw video and emits a ball trajectory, court geometry, and bounce events
raises questions of model coupling, export and packaging, and end-to-end
throughput that do not arise when models are evaluated in isolation. The
purpose of this integrated system is to support automated line calling: by
composing the court geometry with the detected bounce events, the pipeline
determines, for each bounce, whether the ball landed inside or outside the
court boundary. This in/out decision is obtained geometrically, by projecting
the bounce position through the estimated homography into court-metric
coordinates and testing for containment within the court-boundary polygon; it
is qualitative, because no ground-truth line-call labels are available in
either dataset.

\section{Objectives and Research Questions}
\label{sec:intro-objectives}

The overall objective of this thesis is to design, implement, and rigorously
compare representative models for each sub-task on a common protocol, and to
demonstrate that the strongest models can be integrated into a deployable
pipeline that supports automated line calling. Concretely, the work aims to:
(i) construct a deterministic, reproducible benchmark for each sub-task;
(ii) train and evaluate competing models under a shared metric suite;
(iii) integrate the best-performing models into one end-to-end system that
determines, for each detected bounce, whether the ball landed inside or outside
the court (an automated in/out line call); and (iv) release the artifacts
required to reproduce the results. This objective is articulated through four
research questions (RQs).

\begin{description}[leftmargin=2.2em,style=nextline]
  \item[RQ1 (Ball tracking).] Among the TrackNet family (V1--V5) and a
    single-frame detector (YOLO11m), which most accurately localises a
    small, fast, and frequently occluded tennis ball, and what does the
    architectural progression buy in accuracy relative to inference speed?
  \item[RQ2 (Court detection).] Which backbone offers the best accuracy/latency
    trade-off for fourteen-keypoint court detection among a heatmap
    convolutional neural network (CNN), a deep residual network, a
    high-resolution network, and a lightweight mobile network, and is the
    resulting keypoint accuracy sufficient to drive a reliable homography to
    court-metric space?
  \item[RQ3 (Bounce detection).] For bounce-event timing from a noisy ball
    trajectory, how do a hand-built heuristic, a gradient-boosted per-frame
    classifier, and a temporal convolutional network compare in accuracy, and
    what does each cost in engineering effort?
  \item[RQ4 (Integration).] Can the strongest model from each sub-task be
    composed into a single pipeline that ingests raw video and outputs a
    tracked ball, court geometry, bounce events, and, as its culminating
    output, a per-bounce in/out line-call decision obtained by projecting each
    bounce through the court homography, demonstrating the line-calling
    capability as a qualitative geometric result rather than a validated
    accuracy figure, and what throughput and deployment trade-offs (model
    export, per-frame latency, and processing bottlenecks) does such a
    composition incur?
\end{description}

\section{Contributions}
\label{sec:intro-contributions}

In addressing these questions, the thesis makes the following contributions.

\begin{enumerate}[leftmargin=2.2em]
  \item A deterministic, reproducible benchmark protocol for each
    sub-task, comprising a fixed data-split contract and a defined metric suite
    per task (a per-clip split of the TrackNet dataset for ball tracking and
    bounce detection, and a deterministic hash-based split of the court-keypoint
    dataset for court detection), enabling controlled, like-for-like comparison
    within each sub-task.
  \item Building on that protocol, a systematic comparison of six
    ball-tracking, four court-detection, and three bounce-detection models, each
    evaluated on the dataset appropriate to its task, in contrast to the common
    practice of reporting only a single proposed model.
  \item An integrated, packaged pipeline that wires the strongest
    model from each sub-task into one system, with neural-network export to the
    Open Neural Network Exchange (ONNX) format and an interactive demonstration
    interface, illustrating practical deployability. By combining the estimated
    court homography with the detected bounces, the pipeline further produces an
    automated in/out line-call decision for each bounce, demonstrating the
    line-calling application that names the thesis.
  \item Finally, the reproducibility artifacts themselves (split
    manifests, training configurations, metric outputs, and visualisations),
    collected in a single repository.
\end{enumerate}

These contributions are framed as a controlled comparative study on two
public-style datasets rather than as a claim of state-of-the-art performance on
a competition benchmark; the scope and its limitations are stated in
Section~\ref{sec:intro-scope} and revisited in Chapter~\ref{ch:discussion}.

\section{Scope and Limitations}
\label{sec:intro-scope}

The scope of this thesis is deliberately bounded. The study concerns
offline analysis of recorded monocular tennis video; it does not target
live, real-time operation, and the measured per-frame latency of the integrated
pipeline on commodity CPU hardware does not support a real-time claim
(Chapter~\ref{ch:integrated}). The evaluation rests on two datasets (one for
ball and bounce, one for court), so the conclusions pertain to the conditions
those datasets represent and may not transfer without adaptation to markedly
different camera angles, court surfaces, or broadcast styles. The
bounce-detection test set in particular contains relatively few bounce events,
so the corresponding figures should be read with that sample size in mind. The
in/out line-call output is likewise a qualitative geometric demonstration: the
available datasets provide no ground-truth line-call labels, and the
reprojection and bounce-localisation uncertainties are not propagated into a
calibrated confidence, so the in/out determination is not quantitatively
validated and the system is not presented as a calibrated, officiating-grade
line-calling instrument, consistent with the cautionary framing of
Section~\ref{sec:intro-background} \cite{collins2008hawkeye}.
Finally, the system addresses ball, court, and bounce perception only: it does
not perform player tracking, pose estimation, shot classification, multi-camera
fusion, or three-dimensional trajectory reconstruction. These omissions define directions for future work (Chapter~\ref{ch:conclusion}).

\section{Thesis Organisation}
\label{sec:intro-organisation}

The remainder of this thesis is organised as follows.
Chapter~\ref{ch:related-work} reviews background and related work across
small-object detection, ball tracking, court registration, and trajectory-based
event detection, and situates the research gap.
Chapter~\ref{ch:methodology} describes the two datasets, the split policies, and
the shared metric vocabulary.
Chapter~\ref{ch:ball} presents the ball-tracking comparison;
Chapter~\ref{ch:court} the court keypoint detection and homography pipeline; and
Chapter~\ref{ch:bounce} the bounce-event detection comparison.
Chapter~\ref{ch:integrated} describes the integrated system, its model-selection
rationale, and its deployment characteristics.
Chapter~\ref{ch:discussion} discusses cross-cutting findings, threats to
validity, and reproducibility, and Chapter~\ref{ch:conclusion} concludes by
answering the research questions and outlining future work.
